\documentclass[10pt,a4paper]{amsart}
\usepackage[margin=19mm]{geometry}
\usepackage{amsmath,amssymb,mathtools}
\usepackage[scr=rsfs,cal=euler]{mathalfa}
\usepackage{xfrac}
\usepackage[unicode,hidelinks,bookmarks=false]{hyperref}
\newcommand{\ie}{i.e.\ }
\newcommand{\cf}{cf.\ }
\newcommand{\resp}{respectively\ }
\newcommand{\Subsection}{\subsection}
\providecommand{\nobreakdash}{\nobreak\hbox{-}}
\setlength{\emergencystretch}{2em}
\multlinegap0pt
\usepackage{mathtools}
\usepackage{float}
\usepackage{amssymb}
\usepackage{tikz}
\newtheorem{lemma}{Lemma}
\def\E{\boldsymbol{E}}
\newcommand{\GL}{L}
\newcommand{\Om}{\Omega}
\newcommand{\bfx}{\boldsymbol{x}}
\newcommand{\curlvec}{\boldsymbol{\mrm{curl}}\,}
\newcommand{\dsp}{\displaystyle}
\newcommand{\mrm}[1]{\mathrm{#1}}
\title{Trapped modes in electromagnetic waveguides}
\author{Anne-Sophie Bonnet-Ben Dhia, Lucas Chesnel, Sonia Fliss}
\date{Source: arXiv:2512.17763v2 (CC BY 4.0)}
\begin{document}
\maketitle

\noindent\emph{Source and licence.} Trapped modes in electromagnetic waveguides. Version arXiv:2512.17763v2 (CC BY 4.0), distributed by arXiv under the Creative Commons Attribution 4.0 International licence, http://creativecommons.org/licenses/by/4.0/, as recorded in the arXiv licence metadata for this exact version and shown on the abstract page https://arxiv.org/abs/2512.17763v2. Full licence notice: \texttt{licenses/arxiv-2512.17763-CC-BY-4.0.txt}.\par\smallskip

\noindent\textit{Editorial scope.} Counted unit: the article's lemma EstimAbove (source0.tex lines 1151--1156). The article's complete original proof of it is delivered with it (source0.tex lines 1157--1212, one \texttt{proof} environment of 56 lines). No mathematics is written, repaired or re-derived: the statement and the proof are the article's own lines, in the article's own notation, and no retained line is substituted or re-flowed.  No further source-local context is needed: every cross-reference of every retained line resolves inside the two delivered spans.  Nothing was omitted from any retained span, so the package carries no omission record.  No retained line contains a citation, so the wrapper carries no bibliography and no bibliography entry.  No retained line contains a figure environment, an image-inclusion command or drawing code, so no image is delivered and the package declares no asset dependency.  Editorial formatting only, in addition: an AMS \texttt{amsart} preamble that restates the article's own declarations and macros used by the retained lines, namely the environments \texttt{lemma} on separate counters, together with the standard packages those lines need.  The retained environments print in this document as Lemma 1 in order of appearance, whereas the article prints the same items with the numbering of its own full document.  The complete native source of the article as distributed in the arXiv e-print for this exact version is bundled unchanged as \texttt{sources/191343/source0.tex}.
\begin{lemma}
There holds
\begin{equation}\label{EstimAbove}
\dsp\int_{\Om}|\curlvec \tilde{\E}_p|^2\,d\bfx=3\lambda_\bullet\int_\GL\varphi^2\,dxdy-6\int_0^1\varphi(1,y)^2\,dy.
\end{equation}
\end{lemma}

\begin{proof}
In $\Om$, we have $\curlvec \tilde{\E}_p=\curlvec \E_p$. Moreover, in $\Pi_x$ we find
\[
\curlvec \E_p=\curlvec \left(\begin{array}{c}
0\\
E_y\\
E_z
\end{array}\right)=\left(\begin{array}{c}
\partial_yE_z-\partial_zE_y\\
-\partial_xE_z\\
\partial_xE_y
\end{array}\right).
\]
Thus we obtain
\[
\begin{array}{rcl}
\dsp\int_{\Pi_x}|\curlvec \E_p|^2\,d\bfx&=&\dsp\int_{\Pi_x}(\partial_xE_y)^2+(\partial_zE_y)^2+(\partial_xE_z)^2+(\partial_yE_z)^2-2\partial_yE_z\partial_zE_y\,d\bfx.
\end{array}
\]
By observing that 
\[
\dsp\int_{\Pi_x}\partial_yE_z\partial_zE_y\,d\bfx=\int_1^{+\infty}\int_0^1 \partial_yE_z(x,y)\,dy\int_0^1\partial_zE_y(x,z)\,dzdx=0,
\]
we get 
\[
\dsp\int_{\Pi_x}|\curlvec \E_p|^2\,d\bfx=2\int_{\GL_+}|
\nabla\varphi|^2\,dxdy
\]
where $\GL_+\coloneqq(1;+\infty)\times(0;1)$. Similarly, we obtain
\[
\dsp\int_{\Pi_y}|\curlvec \E_p|^2\,d\bfx=\dsp\int_{\Pi_z}|\curlvec \E_p|^2\,d\bfx=2\int_{\GL_+}|
\nabla\varphi|^2\,dxdy.
\]
On the other hand, in $\mathcal{C}=(0;1)^3$, we have
\[
\curlvec \E_p=\left(\begin{array}{c}
\partial_yE_z-\partial_zE_y\\
\partial_zE_x-\partial_xE_z\\
\partial_xE_y-\partial_yE_x
\end{array}\right).
\]
Therefore, we find
\begin{equation}\label{TermLemma}
\dsp\int_{\mathcal{C}}|\curlvec \E_p|^2\,d\bfx=3\int_{\square}|
\nabla\varphi|^2\,dxdy-2\dsp\int_{\mathcal{C}} \partial_yE_z\partial_zE_y+\partial_xE_z\partial_zE_x+\partial_xE_y\partial_yE_x\,d\bfx
\end{equation}
with $\square=(0;1)^2$. But 
\[
\dsp\int_{\mathcal{C}} \partial_yE_z\partial_zE_y\,d\bfx=\int_0^{1}\int_0^1 \partial_yE_z(x,y)\,dy\int_0^1\partial_zE_y(x,z)\,dzdx=\int_0^{1}E_z(x,1)E_y(x,1)\,dx=\int_0^1\varphi(x,1)^2\,dx.
\]
By exploiting that $\varphi(x,y)=\varphi(y,x)$ in $\GL$, similarly we get
\[
\dsp\int_{\mathcal{C}} \partial_xE_z\partial_zE_x\,d\bfx=\dsp\int_{\mathcal{C}} \partial_xE_y\partial_yE_x\,d\bfx=\int_0^1\varphi(x,1)^2\,dx=\int_0^1\varphi(1,y)^2\,dy.
\]
With (\ref{TermLemma}), we deduce (\ref{EstimAbove}).
\end{proof}

\end{document}
